% =============================================================================
%  Phase-I IEEE report - Systolic Array Matrix Multiplication Accelerator
%  Course : CSA-222 Modern Computer Architecture
%  Build  : pdflatex phase1_report.tex  (run twice)   -- or upload to Overleaf
%  Fill in every [PLACEHOLDER] before submitting.
% =============================================================================
\documentclass[conference]{IEEEtran}
\IEEEoverridecommandlockouts
\usepackage{cite}
\usepackage{amsmath,amssymb}
\usepackage{graphicx}
\usepackage{booktabs}
\usepackage{array}
\usepackage{multirow}
\usepackage{xcolor}
\usepackage{url}
\usepackage[hidelinks]{hyperref}
\graphicspath{{../docs/figures/}}

\begin{document}

\title{Design of a 4$\times$4 Output-Stationary Systolic Array Accelerator for Fixed-Point Matrix Multiplication\\
\large Phase I: Architecture, RTL Prototype and Functional Verification}

\author{
\IEEEauthorblockN{Harshvardhan Gupta}
\IEEEauthorblockA{2401010185}
\and
\IEEEauthorblockN{Harsha Gonela}
\IEEEauthorblockA{2401010181}
\and
\IEEEauthorblockN{Abhijeet Dey}
\IEEEauthorblockA{2401010014}
\and
\IEEEauthorblockN{Saumya Kumar}
\IEEEauthorblockA{2401010432}
\and
\IEEEauthorblockN{Shivansh Upadhyay}
\IEEEauthorblockA{2401020109}
\and
\IEEEauthorblockA{\textit{B.Tech CS\&AI / CS\&DS, Semester 5 --- CSA-222 Modern Computer Architecture}\\
\textit{Newton School of Technology}\\
GitHub: \url{https://github.com/Harsh-gupta-07/MCA-Phase1}}
}

\maketitle

% -----------------------------------------------------------------------------
\begin{abstract}
Matrix multiplication dominates the run time of modern machine-learning
workloads, and general-purpose processors execute it inefficiently because
every multiply-accumulate (MAC) operation re-fetches its operands from the
memory hierarchy. Systolic arrays, the structure at the heart of Google's
Tensor Processing Unit, solve this by passing operands rhythmically between
neighbouring processing elements (PEs) so that each value fetched from memory
is reused many times. This paper presents Phase~I of a small,
synthesizable systolic-array accelerator written in Verilog-2005. The design
contains a parameterised $N\times N$ (here $N=4$) output-stationary array of
PEs, each built around a signed Q4.4 fixed-point MAC unit with a 32-bit
accumulator, diagonal input-skew buffers, operand buffers with a simple
matrix data interface, and a five-state controller. A
directed, self-checking verification environment built on Icarus Verilog
and a bit-exact Python golden model checks the MAC unit (10 tests), the PE
(17 checks), the bare array (2 matrix products) and the complete
accelerator (8 matrix products covering identity, signed, fractional and
extreme-value operands). All tests pass, and waveform inspection confirms
the expected diagonal wavefront. One $4\times4$ product completes $3N-2=10$
cycles after the operands start streaming. Throughput analysis, comparison
with a sequential baseline, FPGA synthesis results and architectural
extensions are planned for Phase~II.
\end{abstract}

\begin{IEEEkeywords}
systolic array, matrix multiplication, processing element, multiply-accumulate,
fixed-point arithmetic, hardware accelerator, Verilog
\end{IEEEkeywords}

% -----------------------------------------------------------------------------
\section{Introduction}
For several decades, single-core performance improved ``for free'' as
Dennard scaling let clock frequencies rise without raising power density.
Once that scaling ended, the industry moved first to multicore designs and
then to \emph{domain-specific architectures} (DSAs): hardware tailored to one
class of computation~\cite{hennessy2019golden}. Deep learning is the most
prominent such domain. Its fully connected and convolutional layers reduce
to general matrix multiplication (GEMM), which accounts for most of the
arithmetic in a forward pass~\cite{sze2017efficient}.

A conventional CPU computes $C=A\times B$ with a triple loop. Each MAC needs
two operand loads, so performance is bound by the memory wall rather than
by the arithmetic units. A \emph{systolic array}~\cite{kung1979systolic,kung1982why}
instead arranges a grid of simple PEs that receive operands from their
neighbours, compute, and pass the operands on. One operand fetched from
memory therefore feeds a whole row or column of PEs. The result is high
throughput, short local wires and a regular layout that maps naturally to
silicon or FPGA fabric.

The objective of this project is to design, verify and analyse a small
systolic-array GEMM accelerator in synthesizable Verilog, in line with
Module~6 of the course (heterogeneous computing, TPU/AI hardware) while
reusing the digital-design foundations of Modules~1--2 (adders, registers,
FSMs)~\cite{harris_ddca,patterson_cod}. The project is entirely simulation and synthesis based; no physical
hardware is required. Phase~I, reported here, covers the architecture, the
complete RTL prototype and its functional verification.

% -----------------------------------------------------------------------------
\section{Literature Review}

\subsection{Origins of Systolic Computing}
Kung and Leiserson introduced systolic arrays as a VLSI-friendly way to
implement matrix operations~\cite{kung1979systolic}. Kung later
generalised the idea~\cite{kung1982why}: data flows from memory ``in a
rhythmic fashion, passing through many processing elements before it
returns to memory, much as blood circulates to and from the heart.'' The key
properties are \emph{(i)} high operand reuse, \emph{(ii)} only
nearest-neighbour communication and \emph{(iii)} simple, replicated cells, the same data-parallel principles that underlie SIMD and multiprocessor designs~\cite{culler_pca}.
Together these make the design easy to scale and to pipeline.

\subsection{Systolic Arrays in Modern AI Accelerators}
Google's first TPU~\cite{jouppi2017tpu} uses a $256\times256$ systolic
matrix unit of 8-bit MACs, giving 92~TOPS. The authors report 15--30$\times$
higher inference throughput than contemporary CPUs and GPUs. Later TPU
generations kept the systolic core~\cite{jouppi2021tpuv4i}.
Gemmini~\cite{genc2021gemmini} is an open-source generator of
parameterisable systolic accelerators for RISC-V SoCs, and
SCALE-Sim~\cite{samajdar2018scalesim} is a cycle-accurate simulator used to
explore array dimensions and dataflows. Both show that array size, dataflow
and buffer organisation are first-order design parameters.

\subsection{Dataflow Taxonomy}
Sze \emph{et al.}~\cite{sze2017efficient} classify accelerators by which
operand stays resident in the PE (Table~\ref{tab:dataflow}).
A \emph{weight-stationary} (WS) array, such as the TPU, preloads $B$ into the PEs and
streams $A$ through them, so partial sums move between PEs.
An \emph{output-stationary} (OS) array, such as ShiDianNao~\cite{du2015shidiannao}, keeps each
partial sum $C_{ij}$ in one PE and streams both $A$ and $B$.
Eyeriss~\cite{chen2017eyeriss} proposed a \emph{row-stationary} dataflow
that maximises reuse in convolutions.

\begin{table}[t]
\centering
\caption{Comparison of common systolic dataflows}
\label{tab:dataflow}
\setlength{\tabcolsep}{3pt}
\begin{tabular}{@{}lccc@{}}
\toprule
 & \textbf{Weight-st.} & \textbf{Output-st.} & \textbf{Row-st.}\\
\midrule
Resident in PE & $B$ (weights) & $C$ (psum) & row of filter \\
Moving operands & $A$, psum & $A$, $B$ & mixed \\
Psum movement & PE$\rightarrow$PE & none & local/PE \\
Example & TPU~\cite{jouppi2017tpu} & ShiDianNao~\cite{du2015shidiannao} & Eyeriss~\cite{chen2017eyeriss}\\
Control complexity & medium & low & high \\
\bottomrule
\end{tabular}
\end{table}

We adopt the \textbf{output-stationary} dataflow for three reasons. First,
each accumulator stays inside its PE, so no wide partial sums travel
between PEs and the PE is a direct extension of a single MAC. Second, the
final result is held in the PE and can be read out without a separate
reduction network. Third, the timing relationship is easy to reason about
and to verify, which suits a first prototype. A weight-stationary variant
is planned as a Phase~II comparison point.

\subsection{Fixed-Point Arithmetic for Inference}
Jacob \emph{et al.}~\cite{jacob2018quantization} showed that 8-bit
integer inference loses very little accuracy compared with 32-bit floating point,
while an 8-bit multiplier is roughly an order of magnitude smaller and more
energy-efficient. Production accelerators~\cite{jouppi2017tpu} therefore
use 8-bit operands with wide (32-bit) accumulators. We follow the same
approach with a signed Q4.4 format.

\subsection{Positioning of this Work}
Industrial arrays are too large and too complex to study in a course project.
Our goal is a minimal but complete and fully verified OS array that
exposes the core architectural trade-offs (latency versus throughput, PE
count versus cycles, and buffer and control overhead) at a scale that can be
simulated with open-source tools and synthesised for a mid-range FPGA.

% -----------------------------------------------------------------------------
\section{Proposed Solution}

\subsection{System Architecture}
Fig.~\ref{fig:system} shows the top-level organisation
(\texttt{systolic\_top}). The host (in this phase, the testbench) writes
the operand matrices into the \emph{input buffers} through a simple
write port. After a \texttt{start} pulse, the \emph{controller} streams
one column of $A$ and one row of $B$ per cycle into the \emph{skew
buffers}, which stagger them into the diagonal wavefront the $4\times4$
\emph{PE array} needs. When \texttt{done} is asserted, each element
$C_{ij}$ can be read through the output read port. Table~\ref{tab:modules}
lists the module hierarchy.

\begin{figure*}[t]
\centering
\includegraphics[width=0.92\textwidth]{fig_system_block.png}
\caption{Top-level block diagram of the Phase-I accelerator (\texttt{systolic\_top}).}
\label{fig:system}
\end{figure*}

\begin{table}[t]
\centering
\caption{RTL module hierarchy (Verilog-2005)}
\label{tab:modules}
\setlength{\tabcolsep}{3pt}
\begin{tabular}{@{}lp{4.9cm}@{}}
\toprule
\textbf{Module} & \textbf{Function} \\
\midrule
\texttt{mac\_unit} & signed Q4.4$\times$Q4.4 multiply, 32-bit accumulate \\
\texttt{pe} & MAC + east/south forwarding registers \\
\texttt{systolic\_array} & $N\times N$ PE grid built with \texttt{generate} \\
\texttt{input\_skew} & lane $k$ delayed by $k$ cycles \\
\texttt{matrix\_buffer} & $A$, $B$ storage + write/read interface \\
\texttt{systolic\_controller} & 5-state FSM: sequencing and control \\
\texttt{systolic\_top} & integration + output read mux \\
\bottomrule
\end{tabular}
\end{table}

\subsection{Number Representation}
Operands use signed two's-complement Q4.4: 8~bits with 4 fractional bits,
so the range is $[-8,\,7.9375]$ and the resolution is $0.0625$. The product of two Q4.4
numbers is exactly representable in Q8.8 (16~bits). Products are sign-extended
into a 32-bit Q24.8 accumulator. For $N=4$ the worst-case magnitude is
$4\times128\times128=65\,536$, far below $2^{31}$, so no overflow can occur
inside the array. The real value of a result is its raw code divided by
$2^8$. For example, $1.5\times2.25$ gives codes $\texttt{0x18}\times\texttt{0x24}=864$,
and $864/256=3.375$.

\subsection{MAC Unit and Processing Element}
The \texttt{mac\_unit} (Fig.~\ref{fig:pe}) contains a combinational signed
multiplier, a sign-extension stage, a 32-bit adder and an accumulator
register:
\begin{equation}
\mathrm{acc}_{t+1} =
\begin{cases}
0 & \text{if } \mathit{clear}\\
\mathrm{acc}_t + a_t\, b_t & \text{if } \mathit{en}\\
\mathrm{acc}_t & \text{otherwise}
\end{cases}
\end{equation}
The PE wraps one MAC and adds two forwarding registers. The A operand
received from the west is registered and driven east, and the B operand
received from the north is registered and driven south. Each forwarding
register adds one cycle of latency per hop, which keeps every wire local
and the critical path short. \textit{clear} zeroes both the accumulator and
the forwarding registers, so every computation starts from a clean pipeline.

\begin{figure}[t]
\centering
\includegraphics[width=0.85\columnwidth]{fig_pe.png}
\caption{Processing Element microarchitecture (output-stationary).}
\label{fig:pe}
\end{figure}

\subsection{Array Organisation and Systolic Data Movement}
\texttt{systolic\_array} instantiates $N^2$ PEs in a two-level
\texttt{generate} loop. Row $i$ of $A$ enters PE$(i,0)$ and moves east.
Column $j$ of $B$ enters PE$(0,j)$ and moves south. For element
$a_{ik}$ and element $b_{kj}$ to meet in PE$(i,j)$ in the same cycle, the inputs are
skewed (Fig.~\ref{fig:dataflow}). The controller presents step $k$ (column
$k$ of $A$, row $k$ of $B$) in cycle $k$. \texttt{input\_skew} delays lane
$i$ by $i$ cycles, and the array adds $j$ (or $i$) hops, so
\begin{equation}
t_{\text{meet}}(i,j,k) = k + i + j .
\end{equation}
The last product reaches PE$(N-1,N-1)$ at $t=3N-3$, so a full product
needs $N$ feed cycles plus $2N-2$ drain cycles, i.e.\ $3N-2$ cycles. The skew logic
costs $N(N-1)/2$ byte registers per side. Fig.~\ref{fig:dataflow} also
shows that a PE is idle while the wavefront has not yet reached it or has
already passed it. Phase~II will quantify this PE utilisation and reduce
it by overlapping consecutive matrices.

\begin{figure}[t]
\centering
\includegraphics[width=\columnwidth]{fig_dataflow.png}
\caption{Skewed input wavefront: $A$ rows delayed by $i$, $B$ columns by $j$ cycles.}
\label{fig:dataflow}
\end{figure}

\subsection{Input Buffers and Matrix Data Interface}
\texttt{matrix\_buffer} stores $A$ and $B$ row-major in register
files. The host writes one element per cycle using
(\texttt{wr\_sel}, \texttt{wr\_addr}$=r\cdot N+c$, \texttt{wr\_data}).
On the read side, the controller supplies a step index $k$. The buffer
outputs the vector $\{A_{0k},\dots,A_{N-1,k}\}$ and the vector
$\{B_{k0},\dots,B_{k,N-1}\}$, or zeros when \texttt{feed} is low.
Injecting zeros is harmless to an accumulator and keeps the
drain phase simple. The results remain in the PE accumulators (the output
registers) and are read through a combinational multiplexer addressed by
\texttt{rd\_addr}.

\subsection{Controller State Machine}
The controller (Fig.~\ref{fig:fsm}) is a Moore FSM with five states.
\textsc{Idle} waits for \texttt{start}. \textsc{Clear} zeroes the datapath
for one cycle. \textsc{Feed} lasts $N$ cycles, during which $k$ is
incremented and $\mathit{en}=\mathit{feed}=1$. \textsc{Drain} lasts $2N-2$
cycles with $\mathit{en}=1$ while zeros flush the pipeline. \textsc{Done}
issues a one-cycle \texttt{done} pulse. A \texttt{busy} flag lets the host
avoid overwriting operands while a computation is running.

\begin{figure}[t]
\centering
\includegraphics[width=\columnwidth]{fig_fsm.png}
\caption{State-transition diagram of \texttt{systolic\_controller}.}
\label{fig:fsm}
\end{figure}

\subsection{Verification Methodology}
Verification is bottom-up and every testbench is self-checking: it
prints PASS/FAIL and does not require manual inspection
(Table~\ref{tab:verif}). The reference results for the system-level test
come from \texttt{golden\_model.py}, a pure-Python model that performs
exactly the integer MACs of the hardware. It writes the operands and
expected results to a \texttt{\$readmemh} file. The directed cases are:
$A\times I$, $I\times B$, positive integers, mixed-sign integers, genuine
Q4.4 fractions, all $-8.0$ (largest positive products), $+7.9375\times-8.0$
(largest negative products) and a zero matrix. All RTL is written in Verilog-2005~\cite{ieee1364}, and all simulations use Icarus
Verilog~\cite{icarus} with \texttt{-g2005 -Wall}, and waveforms are inspected in GTKWave.
Yosys~\cite{wolf_yosys} confirms that the RTL synthesises without warnings.

\begin{table}[t]
\centering
\caption{Phase-I verification summary (Icarus Verilog)}
\label{tab:verif}
\setlength{\tabcolsep}{3pt}
\begin{tabular}{@{}llcc@{}}
\toprule
\textbf{Testbench} & \textbf{Checks} & \textbf{\#} & \textbf{Status}\\
\midrule
\texttt{tb\_mac\_unit} & sign, acc., hold, clear, Q4.4, extremes & 10 & Pass \\
\texttt{tb\_pe} & forward, accumulate, freeze, clear & 17 & Pass \\
\texttt{tb\_systolic\_array} & $A\times I$, signed $A\times B$ & 2 & Pass \\
\texttt{tb\_systolic\_top} & 8 golden-model products & 8 & Pass \\
\bottomrule
\end{tabular}
\end{table}

\subsection{Prototype Status and Waveform Evidence}
All four testbenches pass. Fig.~\ref{fig:wave} shows one complete product
(positive-integer test case). The un-skewed vector \texttt{a\_vec}
($A$ column $k$) becomes a diagonal staircase on \texttt{a\_west}, and the
accumulators settle one after another: $C_{00}$ is final after 4 enabled
cycles and $C_{33}$ after 10 ($=3N-2$). The final values
$C_{00}=80$ and $C_{33}=386$ match the golden model. Measured from the
\texttt{start} pulse to the cycle in which the testbench observes
\texttt{done}, a product takes 12 clock cycles: 1~clear, 4~feed, 6~drain
and 1~handshake cycle. This latency figure is preliminary. A detailed
latency and throughput study is deferred to Phase~II.

\begin{figure*}[t]
\centering
\includegraphics[width=0.95\textwidth]{fig_wave_top.png}
\caption{Simulation waveform of one $4\times4$ product (\texttt{tb\_systolic\_top}): FSM states, skewed operand lanes and accumulator values of the diagonal PEs.}
\label{fig:wave}
\end{figure*}

% -----------------------------------------------------------------------------
\section{Plan for Phase II}
Phase~II will turn the functional prototype into a quantitatively evaluated
design:
\begin{enumerate}
\item \textbf{Performance analysis:} hardware cycle and MAC counters;
latency, throughput (MAC/cycle, GOPS at $f_{\max}$) and PE utilisation for
$N\in\{2,3,4,8\}$.
\item \textbf{Baseline comparison:} a sequential single-MAC
multiplier that needs $N^3$ cycles, and the resulting speed-up versus area.
\item \textbf{Throughput optimisations:} double-buffered (ping-pong)
operand buffers and a diagonally propagated clear, so that consecutive
matrices can stream back-to-back.
\item \textbf{Scalability:} tiling of $8\times8$ and $16\times16$
products on the $4\times4$ array.
\item \textbf{Numerics:} rounding and saturation back to Q4.4/Q8.8, with
overflow flags.
\item \textbf{Synthesis:} Vivado synthesis and implementation
for an Artix-7 part, reporting LUT, FF and DSP utilisation and $f_{\max}$
(no physical board required).
\item \textbf{Verification:} constrained-random regression with
$\geq$1000 matrices against the golden model.
\end{enumerate}

% -----------------------------------------------------------------------------
\section{Individual Contributions}
Table~\ref{tab:contrib} lists each member's Phase-I ownership. Every
member authored their own modules, testbench and report section, and
committed them from their own GitHub account.

\begin{table}[h]
\centering
\caption{Individual contributions (Phase I)}
\label{tab:contrib}
\setlength{\tabcolsep}{3pt}
\begin{tabular}{@{}p{1.7cm}p{6.1cm}@{}}
\toprule
\textbf{Member} & \textbf{Contribution} \\
\midrule
{[Member 1]} & Architecture specification and number format; \texttt{mac\_unit.v}, \texttt{tb\_mac\_unit.v}; system block diagram; Abstract and Introduction. \\
{[Member 2]} & \texttt{pe.v}, \texttt{tb\_pe.v}; PE microarchitecture figure; Literature Review (dataflows, fixed-point). \\
{[Member 3]} & \texttt{systolic\_array.v}, \texttt{input\_skew.v}, \texttt{tb\_systolic\_array.v}; dataflow/skew figure and timing analysis. \\
{[Member 4]} & \texttt{matrix\_buffer.v}, \texttt{systolic\_controller.v}, \texttt{systolic\_top.v}; FSM diagram; buffer/controller sections. \\
{[Member 5]} & \texttt{golden\_model.py}, \texttt{tb\_systolic\_top.v}, Makefile and run scripts; waveform figures; verification section, references, report integration and repository maintenance. \\
\bottomrule
\end{tabular}
\end{table}

% -----------------------------------------------------------------------------
\begin{thebibliography}{99}

\bibitem{hennessy2019golden}
J.~L. Hennessy and D.~A. Patterson, ``A new golden age for computer
architecture,'' \emph{Commun. ACM}, vol.~62, no.~2, pp.~48--60, Feb. 2019.

\bibitem{sze2017efficient}
V.~Sze, Y.-H. Chen, T.-J. Yang, and J.~S. Emer, ``Efficient processing of deep
neural networks: A tutorial and survey,'' \emph{Proc. IEEE}, vol.~105, no.~12,
pp.~2295--2329, Dec. 2017.

\bibitem{kung1979systolic}
H.~T. Kung and C.~E. Leiserson, ``Systolic arrays (for VLSI),'' in
\emph{Sparse Matrix Proceedings 1978}, SIAM, 1979, pp.~256--282.

\bibitem{kung1982why}
H.~T. Kung, ``Why systolic architectures?'' \emph{Computer}, vol.~15, no.~1,
pp.~37--46, Jan. 1982.

\bibitem{jouppi2017tpu}
N.~P. Jouppi \emph{et al.}, ``In-datacenter performance analysis of a tensor
processing unit,'' in \emph{Proc. 44th Annu. Int. Symp. Comput. Archit. (ISCA)},
2017, pp.~1--12.

\bibitem{jouppi2021tpuv4i}
N.~P. Jouppi \emph{et al.}, ``Ten lessons from three generations shaped
Google's TPUv4i,'' in \emph{Proc. 48th Annu. Int. Symp. Comput. Archit. (ISCA)},
2021, pp.~1--14.

\bibitem{genc2021gemmini}
H.~Genc \emph{et al.}, ``Gemmini: Enabling systematic deep-learning
architecture evaluation via full-stack integration,'' in \emph{Proc. 58th
ACM/IEEE Design Automation Conf. (DAC)}, 2021, pp.~769--774.

\bibitem{samajdar2018scalesim}
A.~Samajdar, Y.~Zhu, P.~Whatmough, M.~Mattina, and T.~Krishna, ``SCALE-Sim:
Systolic CNN accelerator simulator,'' \emph{arXiv:1811.02883}, 2018.

\bibitem{du2015shidiannao}
Z.~Du \emph{et al.}, ``ShiDianNao: Shifting vision processing closer to the
sensor,'' in \emph{Proc. 42nd Annu. Int. Symp. Comput. Archit. (ISCA)}, 2015,
pp.~92--104.

\bibitem{chen2017eyeriss}
Y.-H. Chen, T.~Krishna, J.~S. Emer, and V.~Sze, ``Eyeriss: An energy-efficient
reconfigurable accelerator for deep convolutional neural networks,''
\emph{IEEE J. Solid-State Circuits}, vol.~52, no.~1, pp.~127--138, Jan. 2017.

\bibitem{jacob2018quantization}
B.~Jacob \emph{et al.}, ``Quantization and training of neural networks for
efficient integer-arithmetic-only inference,'' in \emph{Proc. IEEE/CVF Conf.
Comput. Vis. Pattern Recognit. (CVPR)}, 2018, pp.~2704--2713.

\bibitem{patterson_cod}
D.~A. Patterson and J.~L. Hennessy, \emph{Computer Organization and Design:
The Hardware/Software Interface, RISC-V Edition}, 2nd~ed. Cambridge, MA, USA:
Morgan Kaufmann, 2020.

\bibitem{harris_ddca}
D.~M. Harris and S.~L. Harris, \emph{Digital Design and Computer
Architecture}, 2nd~ed. San Francisco, CA, USA: Morgan Kaufmann, 2012.

\bibitem{culler_pca}
D.~E. Culler, J.~P. Singh, and A.~Gupta, \emph{Parallel Computer
Architecture: A Hardware/Software Approach}. San Francisco, CA, USA: Morgan
Kaufmann, 1999.

\bibitem{ieee1364}
\emph{IEEE Standard for Verilog Hardware Description Language}, IEEE Std
1364-2005, 2006.

\bibitem{icarus}
S.~Williams, ``Icarus Verilog,'' [Online]. Available:
\url{https://steveicarus.github.io/iverilog/}

\bibitem{wolf_yosys}
C.~Wolf, ``Yosys Open SYnthesis Suite,'' [Online]. Available:
\url{https://yosyshq.net/yosys/}

\end{thebibliography}

\end{document}
